% Mechanistic Validity — NEMI 2026 poster, v5
% Five-row layout: one row per validity type.
% Left column: criteria table. Right columns: case study failures.
%
%   cd poster && lualatex mechval_poster_v5.tex

\documentclass[25pt, a0paper, portrait]{tikzposter}

\usepackage[T1]{fontenc}
\usepackage{amsmath,amssymb}
\usepackage{array}
\usepackage{booktabs}
\usepackage{enumitem}
\usepackage{microtype}

\tikzposterlatexaffectionproofoff

% ── Colors: one per validity type ─────────────────────
\definecolor{darktext}{HTML}{2C3E50}
\definecolor{blockbg}{HTML}{FFFFFF}
\definecolor{constructgreen}{HTML}{27AE60}
\definecolor{measureblue}{HTML}{2980B9}
\definecolor{internalorange}{HTML}{E67E22}
\definecolor{externalpurple}{HTML}{8E44AD}
\definecolor{interpteal}{HTML}{16A085}
\definecolor{failred}{HTML}{E74C3C}

% ── Theme ─────────────────────────────────────────────
\usetheme{Default}

\definecolorstyle{MechValColors}{
  \definecolor{colorOne}{HTML}{2C3E50}
  \definecolor{colorTwo}{HTML}{3498DB}
  \definecolor{colorThree}{HTML}{ECF0F1}
}{
  \colorlet{backgroundcolor}{colorThree}
  \colorlet{framecolor}{colorOne}
  \colorlet{titlefgcolor}{white}
  \colorlet{titlebgcolor}{colorOne}
  \colorlet{blocktitlebgcolor}{colorTwo}
  \colorlet{blocktitlefgcolor}{white}
  \colorlet{blockbodybgcolor}{blockbg}
  \colorlet{blockbodyfgcolor}{darktext}
  \colorlet{notefgcolor}{darktext}
  \colorlet{notebgcolor}{colorTwo!12}
  \colorlet{notefrcolor}{colorTwo}
}
\usecolorstyle{MechValColors}

\defineblockstyle{MechValBlock}{
  titlewidthscale=1.0, bodywidthscale=1.0, titleleft,
  titleoffsetx=0pt, titleoffsety=0pt,
  bodyoffsetx=0pt, bodyoffsety=0pt, bodyverticalshift=0pt,
  roundedcorners=12, linewidth=2pt, titleinnersep=10pt, bodyinnersep=12pt
}{
  \draw[rounded corners=\blockroundedcorners, color=blocktitlebgcolor,
        fill=blockbodybgcolor, line width=\blocklinewidth]
    (blockbody.south west) rectangle (blocktitle.north east);
  \ifBlockHasTitle
    \draw[rounded corners=\blockroundedcorners,
          fill=blocktitlebgcolor, color=blocktitlebgcolor,
          line width=\blocklinewidth]
      (blocktitle.south west) rectangle (blocktitle.north east);
  \fi
}
\useblockstyle{MechValBlock}

\definetitlestyle{MechValTitle}{
  width=\paperwidth, roundedcorners=0, linewidth=0pt, innersep=20pt,
  titletotopverticalspace=12mm, titletoblockverticalspace=18mm
}{
  \draw[fill=titlebgcolor, color=titlebgcolor]
    (\titleposleft, \titleposbottom)
    rectangle (\titleposright, \titlepostop);
}
\usetitlestyle{MechValTitle}

\title{%
  \parbox{0.92\linewidth}{\centering\bfseries
    Mechanistic Validity\\[6pt]
    {\large 36 Criteria from 8 Disciplines, Applied to 16 Published Claims}}}
\author{Elliot Tower\quad\texttt{elliot@elliottower.ai}}
\institute{}

\begin{document}

\maketitle[width=\paperwidth]

% ════════════════════════════════════════════════════════
% ROW 1: CONSTRUCT VALIDITY
% ════════════════════════════════════════════════════════
\begin{columns}
\column{0.28}

{
\colorlet{blocktitlebgcolor}{constructgreen}
\block{Construct (C1--C6)}{
  \textit{Is the target concept well-defined?}

  \vspace{3mm}
  A claim is bounded by its weakest validity
  dimension. Five types form a dependency chain:
  construct validity is a ceiling on all others.

  \vspace{4mm}
  \renewcommand{\arraystretch}{1.15}
  \begin{tabular}{@{}cl@{}}
    C1 & Falsifiability \\
    C2 & Structural plausibility \\
    C3 & Convergent validity \\
    C4 & Discriminant validity \\
    C5 & Nomological validity \\
    C6 & Complementation \\
  \end{tabular}

  \vspace{3mm}
  {\footnotesize Philosophy of science, psychometrics,
  genetics.}
}
}

\column{0.36}

{
\colorlet{blocktitlebgcolor}{constructgreen}
\block{C4: Discriminant Validity}{
  \textbf{Does the measure distinguish this construct from
  neighboring constructs?}

  \vspace{3mm}
  \textbf{Knowledge Neurons} (Dai 2022): neurons identified
  by activation magnitude, but the criterion cannot
  separate factual knowledge from linguistic regularity.
  Suppressing these neurons produces off-target effects
  worse than suppressing random neurons.

  \vspace{3mm}
  \textbf{Gender Bias Circuits} (Vig 2020): heads that
  correlate with gendered outputs, but the method never
  tests whether these heads encode bias specifically
  or linguistic competence generally.

  \vspace{3mm}
  {\footnotesize Psychometrics (Campbell \& Fiske 1959).}
}
}

\column{0.36}

{
\colorlet{blocktitlebgcolor}{constructgreen}
\block{C6: Complementation}{
  \textbf{Are labeled subdivisions functionally distinct?}

  \vspace{3mm}
  IOI identifies ``name movers,'' ``S-inhibition heads,'' and
  ``backup name movers.'' Whether these labels name distinct
  functional units or overlapping contributions to one
  mechanism has never been tested.

  \vspace{3mm}
  Benzer's cis-trans test (1955): run two forward passes,
  one with each component intact, and construct a hybrid
  activation. If the hybrid recovers, the components serve
  independent roles. If the hybrid fails, they participate in
  the same functional unit.

  \vspace{3mm}
  {\footnotesize Genetics (Benzer 1955, cis-trans test).}
}
}

\end{columns}

% ════════════════════════════════════════════════════════
% ROW 2: MEASUREMENT VALIDITY
% ════════════════════════════════════════════════════════
\begin{columns}
\column{0.28}

{
\colorlet{blocktitlebgcolor}{measureblue}
\block{Measurement (M1--M7)}{
  \textit{Does the metric measure what it claims?}

  \vspace{4mm}
  \renewcommand{\arraystretch}{1.15}
  \begin{tabular}{@{}cl@{}}
    M1 & Reliability \\
    M2 & Baseline separation \\
    M3 & Stability \\
    M4 & Calibration \\
    M5 & Sensitivity \\
    M6 & Invariance \\
    M7 & Selection correction \\
  \end{tabular}

  \vspace{4mm}
  M4 (calibration) is attempted in every audited
  claim. None confirm it: the numbers are reported
  without establishing what they mean on an
  absolute scale.

  \vspace{3mm}
  {\footnotesize Psychometrics, statistics.}
}
}

\column{0.36}

{
\colorlet{blocktitlebgcolor}{measureblue}
\block{M2: Baseline Separation}{
  \textbf{Is the score distinguishable from a random baseline?}

  \vspace{3mm}
  Three independent failures, all caught by one criterion:

  \vspace{2mm}
  \textbf{SAE interpretability scores} produce the same
  distribution on random, untrained models as on trained ones.
  The metric cannot distinguish learned features from noise.

  \vspace{2mm}
  \textbf{IIA baselines} (Sutter): interchange intervention
  accuracy on some tasks exceeds majority-class rates before
  any causal structure is identified.

  \vspace{2mm}
  \textbf{SAE reconstruction} (Heap): scores that appear strong
  are indistinguishable from a baseline that ignores the learned
  dictionary.

  \vspace{3mm}
  {\footnotesize Psychometrics (floor and ceiling checks).}
}
}

\column{0.36}

{
\colorlet{blocktitlebgcolor}{measureblue}
\block{M6: Measurement Invariance}{
  \textbf{Does the metric behave consistently across conditions?}

  \vspace{3mm}
  \textbf{IOI circuit faithfulness:} the same circuit, measured
  by the same underlying concept, produces scores spanning
  \textbf{below 0\% to over 100\%} across six methodological
  choices.

  \vspace{3mm}
  The headline ``87\% faithfulness'' is a point in a
  distribution that includes contradictory values depending on
  intervention type, metric formula, and ablation target.

  \vspace{3mm}
  If the instrument's reading changes when you swap the
  instrument, the reading is a property of the instrument,
  not the circuit.

  \vspace{3mm}
  {\footnotesize Psychometrics (Meredith's measurement
  invariance).}
}
}

\end{columns}

% ════════════════════════════════════════════════════════
% ROW 3: INTERNAL VALIDITY
% ════════════════════════════════════════════════════════
\begin{columns}
\column{0.28}

{
\colorlet{blocktitlebgcolor}{internalorange}
\block{Internal (I1--I12)}{
  \textit{Is the causal claim warranted?}

  \vspace{4mm}
  \renewcommand{\arraystretch}{1.1}
  {\small
  \begin{tabular}{@{}cl@{}}
    I1 & Necessity \\
    I2 & Sufficiency \\
    I3 & Minimality \\
    I4 & Specificity \\
    I5 & Rival mechanism exclusion \\
    I6 & Double dissociation \\
    I7 & Confound control \\
    I8 & Confounding sensitivity \\
    I9 & Epistatic interaction \\
    I10 & Rescue / reversibility \\
    I11 & Onset coupling \\
    I12 & Offset coupling \\
  \end{tabular}}

  \vspace{3mm}
  {\footnotesize Neuroscience, genetics, causal inference,
  medical microbiology.}
}
}

\column{0.36}

{
\colorlet{blocktitlebgcolor}{internalorange}
\block{I4: Specificity}{
  \textbf{Does the intervention affect this task more than
  matched controls?}

  \vspace{3mm}
  The ORBITA trial (2017): cardiac stenting for stable angina
  produced no benefit over sham surgery. The intervention was
  real; the specificity claim had never been tested against a
  matched control.

  \vspace{3mm}
  \textbf{Knowledge Neurons:} suppressing the identified neurons
  produces off-target effects \textbf{worse than random neurons}.
  The circuit is less specific than chance.

  \vspace{3mm}
  \textbf{IOI:} specificity against a size-matched random
  subnetwork---same count, same layer distribution---has
  never been tested.

  \vspace{3mm}
  {\footnotesize Medicine (sham-controlled trials, ORBITA).}
}
}

\column{0.36}

{
\colorlet{blocktitlebgcolor}{internalorange}
\block{I6: Double Dissociation}{
  \textbf{Do two interventions show crossed effects?}

  \vspace{3mm}
  \renewcommand{\arraystretch}{1.2}
  \begin{tabular}{@{}lcc@{}}
    \toprule
    & \textbf{Behav.\ X} & \textbf{Behav.\ Y} \\
    \midrule
    Ablate mech.\ A & Large deficit & Small deficit \\
    Ablate mech.\ B & Small deficit & Large deficit \\
    \bottomrule
  \end{tabular}

  \vspace{3mm}
  A selective effect in one row is a single dissociation:
  one intervention affects one behavior. The crossed pattern
  across both rows rules out generalized disruption---neither
  effect can be explained by nonspecific damage.

  \vspace{3mm}
  Across audited claims, only induction heads for token
  copying meet this criterion, achieved by Feucht et al.\
  (2025)---three years after the original paper.

  \vspace{3mm}
  {\footnotesize Neuroscience (Shallice 1988).}
}
}

\end{columns}

% ════════════════════════════════════════════════════════
% ROW 4: EXTERNAL VALIDITY
% ════════════════════════════════════════════════════════
\begin{columns}
\column{0.28}

{
\colorlet{blocktitlebgcolor}{externalpurple}
\block{External (E1--E6)}{
  \textit{Does the finding generalize?}

  \vspace{4mm}
  \renewcommand{\arraystretch}{1.15}
  \begin{tabular}{@{}cl@{}}
    E1 & Intervention reach \\
    E2 & Prompt generalization \\
    E3 & Cross-task generalization \\
    E4 & Cross-model generalization \\
    E5 & Graded response \\
    E6 & Novel prediction \\
  \end{tabular}

  \vspace{4mm}
  E2 (prompt generalization) is disconfirmed for
  IOI: the circuit classification changes across
  prompt distributions.

  \vspace{3mm}
  {\footnotesize Causal inference, pharmacology.}
}
}

\column{0.36}

{
\colorlet{blocktitlebgcolor}{externalpurple}
\block{E1: Intervention Reach}{
  \textbf{Do two different intervention methods agree?}

  \vspace{3mm}
  \textbf{IOI} is the sharpest failure. Six published
  evaluations of the same circuit, using different intervention
  methods, produce faithfulness scores spanning \textbf{below
  0\% to over 100\%}.

  \vspace{3mm}
  When independent methods disagree on the \emph{direction} of
  an effect, the result is a property of the method, not the
  circuit.

  \vspace{3mm}
  The field's most studied mechanism, evaluated by the most
  methods, shows the widest disagreement.

  \vspace{3mm}
  {\footnotesize Causal inference (multi-method convergence).}
}
}

\column{0.36}

{
\colorlet{blocktitlebgcolor}{externalpurple}
\block{E5: Graded Response}{
  \textbf{Does partial intervention produce partial effect?}

  \vspace{3mm}
  A dose-response curve between zero and complete ablation
  carries more information than either endpoint. The curve
  shape distinguishes threshold effects from graded mechanisms,
  reveals compensation and saturation, and detects
  non-monotonicity.

  \vspace{3mm}
  Most circuit evaluations report a single ablation point. The
  interpolation between intact and fully ablated is available in
  standard frameworks but almost never recorded.

  \vspace{3mm}
  Genuine mechanisms can produce thresholds; nonspecific
  disruption can produce smooth decline. The curve is
  informative, not dispositive.

  \vspace{3mm}
  {\footnotesize Pharmacology (Hill 1910, dose-response).}
}
}

\end{columns}

% ════════════════════════════════════════════════════════
% ROW 5: INTERPRETIVE VALIDITY + TAKEAWAY
% ════════════════════════════════════════════════════════
\begin{columns}
\column{0.28}

{
\colorlet{blocktitlebgcolor}{interpteal}
\block{Interpretive (V1--V5)}{
  \textit{Is the interpretation justified?}

  \vspace{2mm}
  {\small\textbf{Novel contribution.} No existing validity
  framework includes these criteria.}

  \vspace{4mm}
  \renewcommand{\arraystretch}{1.15}
  \begin{tabular}{@{}cl@{}}
    V1 & Level declaration \\
    V2 & Level-evidence match \\
    V3 & Alternative level \\
    V4 & Unlicensed labeling \\
    V5 & Scope declaration \\
  \end{tabular}

  \vspace{4mm}
  Interpretive validity asks whether the description
  mode matches the evidence. Activation-level evidence
  does not license weight-level claims.

  \vspace{3mm}
  {\footnotesize Marr (1982), Messick (1995).}
}
}

\column{0.36}

{
\colorlet{blocktitlebgcolor}{interpteal}
\block{V4: Unlicensed Labeling}{
  \textbf{Does a name import a property that was not measured?}

  \vspace{3mm}
  \textbf{Othello ``world model''} (Li 2023): a probe recovers
  board state from internal representations. The label ``world
  model'' imports spatial reasoning, planning, and structured
  representation---none of which were contrastively tested.
  Vafa et al.\ later tested counterfactual board states: one
  checkpoint passes, the other fails.

  \vspace{3mm}
  \textbf{Steering vectors} (Arditi 2024): the ``refusal
  direction'' controls refusal behavior across 13 models.
  But controlling a behavior is evidence for a \emph{lever},
  not evidence for how the model \emph{represents} refusal.
  The label imports representational content the evidence
  does not establish.

  \vspace{3mm}
  {\footnotesize Messick (1995, consequential validity).}
}
}

\column{0.36}

{
\colorlet{blocktitlebgcolor}{colorOne}
\block{Takeaway}{
  {\Large\bfseries A mechanistic claim is only as strong as
  its weakest validity dimension.}

  \vspace{4mm}
  \begin{itemize}[leftmargin=*, itemsep=2mm]
    \item \textbf{36} criteria across \textbf{5} validity types
    \item \textbf{8} source disciplines: philosophy of science,
      psychometrics, causal inference, neuroscience, genetics,
      medical microbiology, pharmacology, and Marr/Messick
    \item The field's most studied circuit (IOI) is
      disconfirmed on measurement stability, measurement
      invariance, and prompt generalization
    \item Only one claim reaches Triangulated; none reach
      Validated
    \item Interpretive validity (V1--V5) is a novel
      contribution with no counterpart in existing frameworks
  \end{itemize}

  \vspace{4mm}
  \begin{center}
  \texttt{elliot@elliottower.ai}

  \vspace{2mm}
  {\small Paper:}
  \texttt{doi.org/10.5281/zenodo.20478480}
  \end{center}
}
}

\end{columns}

\end{document}
